% !TEX root = ../main.tex
\chapter{Implementazione}\label{implementazione}

% PATTERN DISPONIBILI IN QUESTO CAPITOLO (i più ricchi della triennale):
%
% LISTATI DI CODICE (linguaggi definiti in Tesi/style/):
%   \begin{lstlisting}[language=yaml, caption={Titolo}, label={lst:nome}]
%   \begin{lstlisting}[language=Rust, caption={Titolo}, label={lst:nome}]
%   \begin{lstlisting}[language=json, caption={Titolo}, label={lst:nome}]
%
% TABELLE CON COLONNE CUSTOM E BOOKTABS:
%   \begin{table}[h!]\centering
%   \begin{tabularx}{\textwidth}{c *{N}{C}}   % C = colonna centrata (tabularx)
%   \toprule
%   Header & Col1 & Col2 \\ \midrule
%   \multirow{2}{*}{testo} & A & B \\          % celle multi-riga verticali
%   \cmidrule(r){2-3}                          % riga orizzontale parziale
%   \bottomrule
%   \end{tabularx}\end{table}
%
% ALBERI DI DIRECTORY:
%   \dirtree{.1 root/.2 subdir/.3 file.json.} % ogni voce termina con punto
%
% ESEMPI FORMALI:
%   \begin{exmp} testo \end{exmp}
%
% FIGURE DA DRAWIO (export PDF):
%   \includegraphics[width=0.9\textwidth]{images/nomefile.drawio.pdf}

\section{}
% TODO
